\chapter{
مقدمه
}
\label{chp:chap1}\thispagestyle{empty}
در این فصل زبان زبان هارا از دیدگاه پارادایم بررسی میکنیم، اصطلاحات و تعاریف مورد نیاز بخش های بعدی را ارایه میکنیم.
% ==========+==========+==========+==========+==========+==========+==========+==========
\section{
مقدمه ای بر زبان های برنامه نویسی
}
واژه پارادایم تعریف دقیق و مشخصی ندارد اما راهی برای درک راحت تر و دسته بندی زبان های برنامه نویسی به دست میدهد. یک پارادایم برنامه نویسی یک شیوه یا نحوه نوشتن برنامه است. در هر زبانی شیوه هایی اسان تر و شیوه هایی سخت تر از دیگران است. 

\subsection{
برخی پارادایم های رایج
}
پارادایم دستوری یا
\eng{imperative}
رایج ترین است. در این نوع برنامه نویسی برنامه به صورت دنباله ای از دستورات صریح نوشته میشود. این دستورات در زمان اجرا تغییراتی در متغیر ها به وجود می آوردند که به عنوان وضعیت یا 
\eng{state}
شناخته میشوند. در زبان های 
\eng{C, Python, MATLAB}
بیشتر از این شیوه استفاده میشود. این زبان ها اصطلاحا 
\eng{Multi-paradigm}
هستند، به این معنی که میتوان در این زبان ها به شیوه های متعدد برنامه نوشت. در زبان های 
\eng{Python, MATLAB}
 پارادایم های شی گرایی و تابعی هم قابل استفاده هستند. در پارادایم شی گرا یا
\eng{object-oriented}
برنامه با تعریف اشیا یا 
\eng{object}
ها که هرکدام وضعیت خود و توابعی که آن وضعیت را تغییر میدهند نوشته میشوند. ایده اصلی روش شی گرا این است که داده ها و توابع با هم جفت باشند. روش تابعی یا
\eng{functional}
بر این ایده مبتنی است که برنامه دیگر وضعیت کلی نداشته باشند و تمام برنامه در چهاچوب توابع نوشته شده باشد. حالت خاص تر پارادایم تابعی، تابعی محض یا
\eng{Purely Functional}
که در آن برنامه به هیچ عنوان وضعیت ندارد، حلقه های تکرار وجود ندارند و توابع میتوانند با هم ترکیب شوند، به صورت ورودی به توابع داده شوند و یا خروجی توابع دیگر باشند. زبان مورد بررسی این گزارش، یعتی زبان
\eng{Haskell}
از دسته زبان های تابعی محض است.

\subsection{
مقایسه دقیق تر پارادایم های تابعی و دستوری
}